(1) 小珠高度为 $h=R(1-\cos\theta)$，距旋转轴的半径为 $R\sin\theta$。惯性系动能为 $T=mR^2[\dot\theta^2+\Omega^2\sin^2\theta]/2$，重力势能 $U=mgR(1-\cos\theta)$。去掉与 $\dot\theta$ 无关的恒定驱动项后，有效势取 $U_{\rm eff}=mgR(1-\cos\theta)-mR^2\Omega^2\sin^2\theta/2$。
(2) 平衡条件为 $U_{\rm eff}'=mR\sin\theta(g-R\Omega^2\cos\theta)=0$。因此总有 $\theta=0$（另有最高点 $\theta=\pi$），并在 $\Omega^2>g/R$ 时出现对称的 $\theta_\pm=\pm\arccos[g/(R\Omega^2)]$。等号时两个新平衡与最低点重合，不是两个互异的新平衡。
(3) 一维约束方程 $mR^2\ddot\theta=-U_{\rm eff}'(\theta)$。最低点二阶曲率 $U_{\rm eff}''(0)=mR(g-R\Omega^2)$，所以 $\omega_{\rm low}^2=g/R-\Omega^2$，当 $\Omega<\sqrt{g/R}$ 时稳定；超过门槛曲率为负，最低点成为不稳定平衡。
(4) 在非零平衡 $\cos\theta_\pm=g/(R\Omega^2)$ 处，二阶曲率为 $mR^2\Omega^2\sin^2\theta_\pm$，故 $\omega_\pm^2=\Omega^2[1-g^2/(R^2\Omega^4)]>0$。临界转速 $\Omega_c=\sqrt{g/R}$ 处最低点刚好失去稳定性，两个对称稳定平衡同时出现，构成理想超临界叉形分岔。
